% Section 3: Problem Formulation

\section{Problem Formulation}

\subsection{Multi-Sensor Optimization Problem}

Following the MUSCAT formulation \cite{muscat2023}, we define the sensor placement problem as a resource allocation optimization. The objective is to determine the optimal number, combination, and placement of sensors in a given Region of Interest (ROI) to maximize system performance under resource constraints. Operationally, UTM / Counter-UAS literature stresses that cooperative (RID/RF-capable) and non-cooperative (dark-target) surveillance are not interchangeable \cite{besada2022cuas,blasch2021information,dulia2024sand}; Section~\ref{sec:related} and our dual-layer objectives encode that split explicitly.

\subsubsection{Sets and Indices}

\begin{itemize}
\item $T$: Set of sensor types, indexed by $t \in T$
\item $L$: Set of potential sensor locations, indexed by $l \in L$
\item $U$: Set of UAV types (characterized by RCS), indexed by $u \in U$
\item $H$: Set of flight altitudes, indexed by $h \in H$
\end{itemize}

\subsubsection{Parameters}

\begin{itemize}
\item $AB$: Available budget (cost limit)
\item $C_t$: Unit cost of sensor type $t$
\item $SC_{tluh}$: Coverage provided by sensor type $t$ at location $l$ detecting UAV $u$ at altitude $h$
\item $\sigma$: Conversion parameter for objective function
\end{itemize}

\subsubsection{Decision Variables}

\begin{equation}
x_{tluh} = \begin{cases}
1 & \text{if sensor } t \text{ is placed at location } l \\
0 & \text{otherwise}
\end{cases}
\end{equation}

\subsubsection{Objective Function}

Maximize total coverage minus cost penalty:

\begin{equation}
\max \sum_{t \in T} \sum_{l \in L} \sum_{u \in U} \sum_{h \in H} SC_{tluh} \cdot x_{tluh} - \sigma \sum_{t \in T} \sum_{l \in L} C_t \cdot x_{tluh}
\label{eq:objective}
\end{equation}

\subsubsection{Constraints}

\begin{equation}
\sum_{t \in T} \sum_{l \in L} \sum_{u \in U} \sum_{h \in H} C_t \cdot x_{tluh} \leq AB
\label{eq:budget}
\end{equation}

\begin{equation}
x_{tluh} \in \{0, 1\} \quad \forall t \in T, l \in L, u \in U, h \in H
\label{eq:binary}
\end{equation}

\subsection{Genetic Algorithm Reformulation}

The binary allocation $x_{tluh}$ above is the MUSCAT-style decision variable \cite{muscat2023}. For GA / NSGA-II search we encode deployments as chromosomes (variable-length lists of activated site--type genes), which is an implementation encoding of the same allocation---not a different physical problem:

\begin{equation}
\textbf{X} = [(l_1, t_1, o_1, a_1), \ldots, (l_{n}, t_{n}, o_{n}, a_{n})]
\label{eq:chromosome}
\end{equation}

where each gene represents:
\begin{itemize}
\item $l_i \in L$: candidate site index (several genes may share a site)
\item $t_i \in T$: sensor type
\item $o_i$: azimuth index in \(\{0^{\circ},120^{\circ},240^{\circ}\}\)
\item $a_i \in \{0, 1\}$: activation bit
\end{itemize}
The chromosome length \(n\) is the maximum sensor count; at least the minimum count of genes is kept active.

The search is not scalarized. Each chromosome is scored by the three dual-layer objectives
\begin{equation}
F(\textbf{X}) = \bigl(M_{wp,\mathrm{coop}}(\textbf{X}),\; M_{wp,\mathrm{noncoop}}(\textbf{X}),\; C(\textbf{X})\bigr)
\label{eq:fitness}
\end{equation}
(maximize, maximize, minimize), defined in Section~\ref{sec:dual_metrics}. CapEx counts hardware plus a site-activation fee for every distinct occupied site:
\begin{equation}
C(\textbf{X}) = \sum_{i:\,a_i=1} C_{t_i} + C_{\mathrm{site}}\,\bigl|\{l_i : a_i = 1\}\bigr|,
\label{eq:capex}
\end{equation}
with \(C_{\mathrm{site}}=\$15\mathrm{k}\) unless stated otherwise.

\subsection{MUSCAT Metrics}

To maintain compatibility with the original framework, we also report the classic metrics below; Controls~A--B use \(M_c\), while all dual-layer fronts use \(M_{wp}\).

\textbf{Coverage Index} (Equation 6 from \cite{muscat2023}):
\begin{equation}
\mc = \frac{n}{N}
\label{eq:coverage_index}
\end{equation}
where $n$ is the number of grid cells with detection probability $P_D \geq$ threshold and $N$ is the total number of cells.

\textbf{Gap Index} (Equation 7 from \cite{muscat2023}):
\begin{equation}
\mg = 1 - \mc
\label{eq:gap_index}
\end{equation}

\textbf{Cost-Effectiveness} (Equation 8 from \cite{muscat2023}):
\begin{equation}
CA = \frac{\sum_{t \in T} C_t \cdot x_t}{\mc}
\label{eq:cost_effectiveness}
\end{equation}

\subsection{Requirements}

Legacy MUSCAT corridor thresholds used high single-score coverage and overlap floors \cite{muscat2023}. In the dual-layer revision we instead post-filter Pareto members with joint planning floors on $M_{wp,\mathrm{coop}}$ and $M_{wp,\mathrm{noncoop}}$ (e.g., $0.90/0.35$), reflecting noncollaborative UTM surveillance needs \cite{besada2022cuas} rather than a fused GREEN label alone.


